%% Recursive formulas related to
%% the summation of the M\"obius function,
%% by
%% M. Benito and J. L. Varona
%%
%%     Contact address: jvarona@dmc.unirioja.es
%%
%% The Open Mathematics Journal, 2008, Vol. 1, 25--34
%%
%% http://www.bentham.org/open/tomatj/
%%
% Last modification of this document: September 5, 2008

\documentclass[10pt]{amsart}
\addtolength{\textwidth}{1.4cm}
\hoffset=-0.7cm
%\addtolength{\textheight}{0.8cm}
\voffset=0.3cm

% With these commands we add 1cm to the separation of the headings:
\addtolength{\topmargin}{-1cm}
\addtolength{\headsep}{1cm}

\usepackage{amsmath,amsthm}
\allowdisplaybreaks

%\date{September 5, 2008}

\newcommand{\lcm}{\operatorname{lcm}}

%\theoremstyle{plain}
\newtheorem{theor}{Theorem}
\newtheorem{propo}[theor]{Proposition}

% The LaTeX Companion 2nd ed., section 4.4.2, p. 242 and ff.
\usepackage{fancyhdr}

\setcounter{page}{25}

\fancypagestyle{firstpage}{%
\fancyhead{} % Clear header fields
\fancyheadoffset[R,L]{1.25cm}
\fancyhead[L]{} % left
\fancyhead[C]{\bfseries\itshape\scriptsize 
  The Open Mathematics Journal, 2008, 1, 25--34} % center
\fancyhead[R]{\bfseries\itshape\scriptsize\thepage} % right
\fancyfoot[C]{\vspace*{0.5cm}\bfseries\itshape\scriptsize 
  1874-1177/08\qquad 2008 Bentham Open} % center
}


\fancypagestyle{usual}{%
\fancyhead{} % Clear header fields
\fancyheadoffset[R,L]{1.25cm}
\fancyhead[LO]{\bfseries\itshape\scriptsize 
  Recursive Formulas Related to the Summation} % left odd
\fancyhead[RO]{\bfseries\itshape\scriptsize
  The Open Mathematics Journal, 2008, Volume 1\qquad\thepage} % right odd
\fancyhead[LE]{\bfseries\itshape \scriptsize
  \thepage\qquad 
  The Open Mathematics Journal, 2008, Volume 1} % left even
\fancyhead[RE]{\bfseries\itshape\scriptsize 
  Benito and Varona} % right even
\fancyfoot{}
\renewcommand{\headrulewidth}{0pt}%
}

\fancypagestyle{final}{%
\fancyfoot[L]{\vspace*{-3cm}%
\hspace*{-0.02\textwidth}\rule{1.04\textwidth}{0.6pt}\hspace*{-0.02\textwidth}\hss\\[1mm]
\sffamily 
Received: January 24, 2008
\hfill
Revised: April 18, 2008
\hfill
Accepted: September 6, 2008\\[3mm]
\copyright\ Benito and Varona; Licensee Bentham Open.\\[3mm]
This is an open access article licensed under the terms of the Creative Commons Attribution Non-Commertial License (http://creativecommons.org/licenses/by-nc/3.0/) which permits unrestricted, non-commertial use, distribution and reproduction in any medium, provided the work is properly cited.}%
}


%%%%%%%%%%%%%%%%%%%%%%%%%%%%%%%%%%%%%%%%%%%%
\begin{document}
%%%%%%%%%%%%%%%%%%%%%%%%%%%%%%%%%%%%%%%%%%%%

%%%%%%%%%%%%%%%%%%%%%%%%%%%%%%%%%%%%%%%%%%%%%%
\title[The summation of the M\"obius function]{%
\vspace*{-10mm}
\hspace*{-0.12\textwidth}\rule{1.24\textwidth}{0.6pt}\hspace*{-0.12\textwidth}\\[1mm]
Recursive formulas related to the summation\\of the M\"obius function}
%%%%%%%%%%%%%%%%%%%%%%%%%%%%%%%%%%%%%%%%%%%%%%

%    Information for first author
\author{Manuel Benito}
%    Address of record for the research reported here
\address{Instituto Sagasta,
Glorieta del Doctor Zub\'{\i}a s/n,
26003 Logro\~no, Spain}
%    Current address
%\curraddr{---}
\email{mbenit8@palmera.pntic.mec.es}

%    Information for second author
\author{Juan L. Varona}
\address{Departamento de Matem\'aticas y Computaci\'on, Universidad de La
Rioja, Edificio J.~L.~Vives, Calle Luis de Ulloa s/n, 26004 Logro\~no, Spain}
\email{jvarona@dmc.unirioja.es}
\urladdr{http://www.unirioja.es/dptos/dmc/jvarona/welcome.html}
%    \thanks will become a 1st page footnote.
\thanks{Research of the second author supported by grant MTM2006-13000-C03-03
of the DGI}

%    General info
\subjclass[2000]{Primary 11A25}

\keywords{M\"obius function, summation of the M\"obius function,
Mertens function}

%%%%%%%%%%%%%%%%%%%%%%%%%%%%%%%%%%%%%%%%%%%%%%
\begin{abstract}
For positive integers $n$, let $\mu(n)$ be the M\"obius function,
and $M(n)$ its sum $M(n) = \sum_{k=1}^n \mu(k)$.
We find some identities and recursive formulas for computing $M(n)$;
in particular, we present a two-parametric family of recursive formulas.
\end{abstract}
%%%%%%%%%%%%%%%%%%%%%%%%%%%%%%%%%%%%%%%%%%%%%%


%%%%%%%%%%%%%%%%%%%%%%%%%%%%%%%%%%%%%%%%%%%%%%
\maketitle
%\thispagestyle{empty}
%\pagestyle{empty}

\thispagestyle{firstpage}
\pagestyle{usual}

%%%%%%%%%%%%%%%%%%%%%%%%%%%%%%%%%%%%%%%%%%%%%%

%%%%%%%%%%%%%%%%%%%%%%%%%%%%%%%%%%%%%%%%%%%%%%
\section{Introduction}
%%%%%%%%%%%%%%%%%%%%%%%%%%%%%%%%%%%%%%%%%%%%%%

The well-known M\"obius function $\mu(n)$ is defined, for positive integers $n$, as
$$
\mu(n) := \left\{
\begin{array}{ll}
1 & \mbox{ if } n=1,\\
(-1)^k & \mbox{ if $n$ is a product of $k$ different prime numbers,}\\
0 & \mbox{ if there exists a prime $p$ such that $p^2$ divides $n$}
\end{array}\right.
$$
(see~\cite[Chapter~2]{Ap}).
Then, for every real number $x\geq 0$, the summation of the M\"obius function
is defined by taking
\begin{equation*}
%\label{eq:M}
M(x) = M(\lfloor x \rfloor) := \sum_{k=1}^{\lfloor x \rfloor}\mu(k).
\end{equation*}
In what follows, and as usually, we refer to $M(x)$ as the Mertens function,
although, before F.~Mertens (who used it in 1897, see~\cite{Mer}), 
T.~J.~Stieltjes already had introduced this function in his attempts to prove the Riemann Hypothesis
(see~\cite[Lettre 79, p.~160--164]{Her}, dated in 1885).

The behaviour of $M(x)$ is rather erratic and difficult of analyze,
but it is very important in analytic number theory.
In 1912, J.~E.~Littlewood \cite{Lit} proved that
the Riemann Hypothesis is equivalent to this fact:
\begin{equation}
\label{eq:Me}
|M(x)| = O(x^{1/2+\varepsilon}), \quad\mbox{when } x\to \infty,
\quad\mbox{for every } \varepsilon>0;
\end{equation}
in relation to this subject, see also~\cite{Tit}. Of course,
it is not yet known if~(\ref{eq:Me}) is true or false.
%
Previously, in 1897, Mertens \cite{Mer} had given a table of values
of $M(n)$ for $1\leq n\leq 10000$. Relying on this table, he conjectured that,
for $x>1$,
$$
|M(x)| < \sqrt{x}.
$$
This conjecture was disproved, in 1985, by A.~M.~Odlyzko and H.~te Riele
\cite{OdRie}, but they did not find an explicit counterexample. Actually,
for every value of $M(n)$ computed up to that date, always happened
$|M(n)| < 0.6 \sqrt n$.
In 1987, J. Pintz \cite{Pi} proved that the Mertens conjecture is false for some 
$n < \exp(3.21 \times 10^{64})$;
and this was improved further recently in 2006 by T.~Kotnik and H.~te Riele \cite{KoRie},
who showed that the Mertens conjecture is false for some $n < \exp(1.59 \times 10^{40})$.
More studies about the order of the Mertens function can be found in~\cite{Ma} and~\cite{KoLu}.
Nowadays, to find an explicit counterexample of the Mertens conjecture
is yet a very pursued result in number theory, and it generally believed
that no counterexample will be found for $n < 10^{20}$.

To evaluate $M(n)$, a big quantity of recursive formulas appear in the mathematical literature.
For instance, Stieltjes~\cite[Letter~79, p.~163]{Her} proved the expression
\begin{equation} 
\label{eq:Stiel}
\sum_{k \le \sqrt n} (-1)^{k-1} M(n/k) 
= -1 + M(\sqrt n) z(\sqrt n) - \sum_{k \le \sqrt n } z(n/k)  \mu(k)
\end{equation}
where $z(x) = 0$ if $\lfloor x\rfloor$ is even and $1$ if it is odd;
some other recursive formulas appear in the famous \textit{Primzahlen} of E.~Landau~\cite{Lan}.
In 1996, M.~Del\'eglise and J.~Rivat \cite{DeRi},
used an algorithm derived from the recurrence formula
\begin{equation}
\label{eq:DR}
M(x) = M(u) - \sum_{a \leq u} \mu(a) \sum_{\frac ua < b \leq \frac xa}
M\left( \frac x{ab}\right)
\end{equation}
(being  $1 \leq u \leq x$) to evaluate $M(10^{16}) = -3195437$.
More recursive formulas can be found in~\cite{Neu}, and~\cite{Dr}; also,
a large number of further references to related studies, including a nice historical review, 
are given in~\cite{Rie}.

The aim of this paper is to prove different identities and recursive
formulas satisfied by the Mertens function~$M$.
We devote to this end sections~\ref{sec:M}, \ref{sec:mu} and~\ref{sec:Mmu};
see Theorems~\ref{th:Ma}, \ref{th:Mk},
\ref{th:2M+3}, \ref{th:2M+3gh}, and~\ref{th:2M+3sr}.
For instance, in Theorem~\ref{th:Mk} we present a formula to
evaluate $M(n)$ similar to the one given by its definition,
but with only $\left\lfloor\frac n3\right\rfloor$ summands.
Also, let us note the interesting expansion for $2M(n)+3$
that appears in Theorem~\ref{th:2M+3}, as well as the properties of the
involved coefficients, studied below; they will lead us to
Theorems~\ref{th:2M+3gh} and~\ref{th:2M+3sr}.
In particular, Theorem~\ref{th:2M+3sr} gives a two-parametric family
of recursive formulas for computing the Mertens function.
As long as we know, all the ``theorems'' that we present in these sections
are new; however, some of the ``propositions'' are already known, and we have included
them by completeness.

Finally, in section~\ref{sec:H}, we study some properties
of a function (that we will denote $H(n,m)$)
related with the ones that appear in the previous
sections; in particular, we prove the periodicity of this function.


%%%%%%%%%%%%%%%%%%%%%%%%%%%%%%%%%%%%%%%%%%%%%%
\section{Formulas in which only $M$ appears}
\label{sec:M}
%%%%%%%%%%%%%%%%%%%%%%%%%%%%%%%%%%%%%%%%%%%%%%

Let us begin by recalling the following well-known property of the
M\"obius function:
\begin{equation}
\label{eq:summu}
\sum_{d|n}\mu(d) =
\left\{
\begin{array}{lll}
1& \mbox{ if } & n=1,\\
0& \mbox{ if } & n>1.
\end{array}
\right.
\end{equation}
Indeed, it is trivial for $n=1$. And, for $n>1$, if
$n = \prod_{j=1}^kp_j^{\alpha_j}>1$ ($p_j$ primes, $p_j \ne p_i$ for $j \ne i$),
then
$$
\sum_{d|n}\mu(d) = \binom{k}{0} - \binom{k}{1}
+ \cdots + (-1)^k\binom{k}{k} = (1-1)^k = 0.
$$
The identity~(\ref{eq:summu}) allows to find a way of relating the value
$M(n)$ with the values of $M(m)$, with $m$ less than~$n$. This result, also
known (and whose proof we reproduce by completeness), is the following:

\begin{propo}
\label{th:1=}
For every positive $n$, the Mertens function verifies
\begin{equation}
\label{eq:1=}
1 = \sum_{a=1}^n M\left(\frac na\right).
\end{equation}
\end{propo}

\begin{proof}
Actually, we will prove~(\ref{eq:1=}) also for real numbers $x \geq 1$.
{}From the definition $M(x) = \sum_{k \leq x}\mu(k)$, we have
$$
\sum_{a=1}^{\lfloor x\rfloor} M\left(\frac xa\right)
= \sum_{a=1}^{\lfloor x\rfloor}
\sum_{b=1}^{\lfloor \frac xa \rfloor} \mu(b).
$$
If $ab=k$, then $a|k$ and, moreover, when the values of $a$ and $b$ vary,
$k$ takes the values $1, 2, \ldots, \lfloor x\rfloor$.
Then, we have
$$
\sum_{a=1}^{\lfloor x\rfloor}
\sum_{b=1}^{\lfloor \frac xa \rfloor} \mu(b)
= \sum_{1\leq k\leq \lfloor x\rfloor}\sum_{a|k}\mu(a).
$$
By applying~(\ref{eq:summu}), we get~(\ref{eq:1=}).
\end{proof}

Of course, from (\ref{eq:1=}) we obtain the following recursive formula
satisfied by~$M(n)$:
\begin{equation}
\label{eq:1-}
M(n) = 1-\sum_{a=2}^n M\left(\frac na\right),
\end{equation}
which is essentially one of the recursive formulae used by Neubauer~\cite{Neu} 
to compute $M(n)$ up to $10^{10}$.
Moreover, let us note that (\ref{eq:summu}) and (\ref{eq:1=}) were used by Del\'eglise
and Rivat~\cite{DeRi} to find the identity~(\ref{eq:DR}).

In (\ref{eq:1-}), $n$ summands appear. In the following
theorem, we reduce the number of summands up to
$\left\lfloor \frac {n-1}2\right\rfloor$.

\begin{theor}
\label{th:Ma}
If $n\geq 3$, then
\begin{equation}
\label{eq:Ma}
M(n) = -\sum_{a=1}^{\lfloor\frac{n-1}2\rfloor} M\left(\frac {n}{2a+1}\right).
\end{equation}
\end{theor}

\begin{proof}
If $n=2m$ with $m>1$, by applying~(\ref{eq:1-}) and~(\ref{eq:1=}),
we get
\begin{align*}
M(2m) &= 1-\sum_{a=2}^{2m} M\left(\frac {2m}{a}\right)
= \sum_{a=1}^{m} M\left(\frac {m}{a}\right)-\sum_{a=2}^{2m} M\left(\frac {2m}{a}\right) \\
&= -\sum_{a=1}^{m-1} M\left(\frac {2m}{2a+1}\right)
= -\sum_{a=1}^{\lfloor\frac{n-1}2\rfloor} M\left(\frac {n}{2a+1}\right).
\end{align*}

For the case $n=2m+1$, let us first note that the greatest remainder 
that can be obtained when $m$ is divided by $a$ is $a-1$, and, moreover
$$
\frac{a-1}a + \frac 1{2a} = \frac{2a-2+1}{2a} = \frac{2a-1}{2a} < 1.
$$
Thus, it is clear that
$$
M\left(\frac{2m+1}{2a}\right)
= M\left(\frac ma + \frac 1{2a}\right)
= M\left(\frac{m}{a}\right).
$$
Then, by applying~(\ref{eq:1-}), (\ref{eq:1=}), and this fact, we get
\begin{align*}
M(2m+1) &= 1-\sum_{a=2}^{2m+1} M\left(\frac {2m+1}{a}\right) %\\
= \sum_{a=1}^{m} M\left(\frac {m}{a}\right)
- \sum_{a=2}^{2m+1} M\left(\frac {2m+1}{a}\right) \\
&= -\sum_{a=1}^{m} M\left(\frac {2m+1}{2a+1}\right)
= -\sum_{a=1}^{\lfloor\frac{n-1}2\rfloor} M\left(\frac {n}{2a+1}\right).
\qedhere
\end{align*}
\end{proof}

%%%%%%%%%%%%%%%%%%%%%%%%%%%%%%%%%%%%%%%%%%%%%
\section{Formulas in which only $\mu$ appears}
\label{sec:mu}
%%%%%%%%%%%%%%%%%%%%%%%%%%%%%%%%%%%%%%%%%%%%%

In the following theorem, we expand $M(n)$ as a sum with
$\left\lfloor\frac{n}3\right\rfloor$
summands, in which only $\mu$ and the integer-part function appear.
In particular, this result provides a more efficient way to
compute $M(n)$ than just to use its definition
$M(n) = \sum_{k=1}^n \mu(k)$.

\begin{theor}
\label{th:Mk}
If $n\geq 3$, then
\begin{equation}
\label{eq:Mk}
M(n) = -\sum_{k=1}^{\lfloor\frac{n}3\rfloor}
\left\lfloor \frac{n-k}{2k}\right\rfloor \mu(k).
\end{equation}
\end{theor}

\begin{proof}
Let us remind~(\ref{eq:Ma}) in Theorem~\ref{th:Ma}.
The greatest value achieved by $\left\lfloor \frac{n}{2a+1} \right\rfloor$ is
$\left\lfloor \frac{n}{3} \right\rfloor$. Moreover, 
$\left\lfloor \frac{n}{2a+1} \right\rfloor$ takes value $k$ if
$$
k \leq \frac{n}{2a+1} < k+1,
$$
i.e.,
$$
\frac{n-(k+1)}{2(k+1)} < a \leq \frac {n-k}{2k}.
$$
In this way, $\left\lfloor \frac{n}{2a+1}\right\rfloor = k$ for
$\left\lfloor \frac{n-k}{2k}\right\rfloor
- \left\lfloor \frac{n-(k+1)}{2(k+1)}\right\rfloor$ values
of~$a$.

As a consequence,
\begin{multline*}
M(n) = -\sum_{k=1}^{\lfloor\frac{n}3\rfloor}
\left( \left\lfloor \frac{n-k}{2k}\right\rfloor
- \left\lfloor \frac {n-(k+1)}{2(k+1)}\right\rfloor \right) M(k)
\\
= -\Bigg(
\left( \left\lfloor \frac{n-1}{2}\right\rfloor
- \left\lfloor \frac {n-2}{2\cdot 2}\right\rfloor \right) M(1)
+ \left( \left\lfloor \frac {n-2}{2\cdot 2}\right\rfloor
- \left\lfloor \frac {n-3}{2\cdot 3}\right\rfloor \right) M(2)
+ \cdots
\\
+ \left( \left\lfloor \frac {n-\left\lfloor \frac n3\right\rfloor}{2\left\lfloor
\frac n3\right\rfloor}\right\rfloor 
- \left\lfloor \frac {n-\left\lfloor \frac n3\right\rfloor-1}{2\left(\left\lfloor
\frac n3\right\rfloor+1\right)}\right\rfloor\right)
M\left(\left\lfloor \frac n3\right\rfloor\right)
\Bigg)
\\
= -\Bigg(
\left\lfloor \frac {n-1}{2}\right\rfloor \mu(1)
+ \left\lfloor \frac {n-2}{2\cdot 2}\right\rfloor \mu(2)
+ \left\lfloor \frac {n-3}{2\cdot 3}\right\rfloor \mu(3)
+ \cdots
\\
+ \left\lfloor \frac {n-\left\lfloor \frac n3\right\rfloor}{2\left\lfloor
\frac n3\right\rfloor}\right\rfloor
\mu\left(\left\lfloor \frac n3\right\rfloor\right)
- \left\lfloor \frac {n-\left\lfloor \frac n3\right\rfloor-1}{2\left(\left\lfloor
\frac n3\right\rfloor+1\right)}\right\rfloor
M\left(\left\lfloor \frac n3\right\rfloor\right)
\Bigg).
\end{multline*}
Now, let us observe
$$
\frac {n-\left\lfloor \frac n3\right\rfloor-1}{2\left(\left\lfloor
\frac n3\right\rfloor+1\right)}
= \left\{ \begin{array}{lll}
\frac{3m-m-1}{2m+2} = \frac{2m-1}{2m+2} < 1 & \mbox{ if } & n=3m ,\\[1mm]
\frac{3m+1-m-1}{2m+2} = \frac{2m}{2m+2} < 1  & \mbox{ if } & n=3m+1,\\[1mm]
\frac{3m+2-m-1}{2m+2} = \frac{2m+1}{2m+2} < 1 & \mbox{ if } & n=3m+2,
\end{array}
\right.
$$
and so
$$
\left\lfloor \frac {n-\left\lfloor \frac n3\right\rfloor-1}{2\left(\left\lfloor
\frac n3\right\rfloor+1\right)}\right\rfloor = 0.
$$
Then, (\ref{eq:Mk}) follows.
\end{proof}

The following result relates the value of $\mu(n)$ to
the values of $\mu(m)$ for $1\leq m<n$. Actually, this result is
already known (see~\cite[Theorem~3.12]{Ap}),
although the proof that we make in this paper is different and,
perhaps, new; here, we use an argument similar to the one used in the proof
of Theorem~\ref{th:Mk}.

\begin{propo}
\label{th:1=k}
The M\"obius function satisfies
\begin{equation}
\label{eq:1=k}
1 = \sum_{k=1}^n \left\lfloor \frac {n}{k}\right\rfloor \mu(k).
\end{equation}
\end{propo}

\begin{proof}
By Proposition~\ref{th:1=}, $1 = \sum_{a=1}^n M\left(\frac na\right)$.
Here, we have $\left\lfloor \frac {n}{a}\right\rfloor =k$ if and only if
$$
k\leq \frac na < k+1,
$$
i.e.,
$$
\frac n{k+1} < a \leq \frac nk,
$$
and so $M\left(\frac na\right) = M\left(\left\lfloor\frac n a\right\rfloor\right) = M(k)$
for $\left\lfloor\frac n{k}\right\rfloor -\left\lfloor\frac n{k+1}\right\rfloor$
values of~$a$.

Then
\begin{align*}
1 &= \sum_{a=1}^n M\left(\frac na\right)
= \sum_{k=1}^n \left( \left\lfloor\frac n{k}\right\rfloor
- \left\lfloor\frac n{k+1}\right\rfloor\right) M(k)
\\
&= \left( \left\lfloor\frac n{1}\right\rfloor -\left\lfloor\frac n{2}\right\rfloor\right) M(1)
+ \left( \left\lfloor\frac n{2}\right\rfloor -\left\lfloor\frac n{3}\right\rfloor\right) M(2)
+ \cdots
+ \left( \left\lfloor\frac n{n}\right\rfloor -\left\lfloor\frac n{n+1}\right\rfloor\right) M(n)
\\
&= \left\lfloor\frac n{1}\right\rfloor
+ \left\lfloor\frac n{2}\right\rfloor (M(2)-M(1))
+ \cdots
+ \left\lfloor\frac n{n}\right\rfloor (M(n)-M(n-1))
- \left\lfloor\frac n{n+1}\right\rfloor M(n)
\\
&= \left\lfloor\frac n{1}\right\rfloor \mu(1)
+ \left\lfloor\frac n{2}\right\rfloor \mu(2)
+ \cdots
+ \left\lfloor\frac n{n}\right\rfloor \mu(n)
= \sum_{k=1}^n \left\lfloor \frac {n}{k}\right\rfloor \mu(k).
\qedhere
\end{align*}
\end{proof}

Now, let us prove another expansion of $1$ as a sum of $\mu$'s.
We will use this result in the proof of Theorem~\ref{th:2M+3}.

\begin{propo}
\label{th:1=3k}
For $n \geq 3$, the M\"obius function satisfies
$$
1 = \sum_{k=1}^{\lfloor\frac{n}3\rfloor}
\left\lfloor \frac {n}{3k}\right\rfloor \mu(k).
$$
\end{propo}

\begin{proof}
First, let us suppose $n=3m$. By Proposition~\ref{th:1=k}, we have
$$
1 = \sum_{k=1}^{m} \left\lfloor \frac{m}{k}\right\rfloor \mu(k)
= \sum_{k=1}^{\frac{n}3} \left\lfloor \frac{3m}{3k}\right\rfloor \mu(k)
= \sum_{k=1}^{\lfloor\frac{n}3\rfloor} \left\lfloor \frac{n}{3k}\right\rfloor \mu(k).
$$

If $n=3m+1$,
$$
\sum_{k=1}^{\lfloor\frac{n}3\rfloor} \left\lfloor \frac {n}{3k}\right\rfloor \mu(k)
= \sum_{k=1}^{m} \left\lfloor \frac{3m+1}{3k}\right\rfloor \mu(k)
= \sum_{k=1}^{m} \left\lfloor \frac{m}{k}\right\rfloor \mu(k) = 1,
$$
because, in
$$
\left\lfloor \frac{3m+1}{3k}\right\rfloor
= \left\lfloor \frac{m}{k} + \frac 1{3k}\right\rfloor,
$$
the remainder when $m$ is divided by $k$ is always less or equal than $k-1$;
and, by being
$$
\frac {k-1}k + \frac 1{3k} = \frac{3k-2}{3k} < 1,
$$
we have
$$
\left\lfloor \frac{3m+1}{3k}\right\rfloor
= \left\lfloor \frac{m}{k}\right\rfloor.
$$

If $n=3m+2$,
$$
\left\lfloor \frac{3m+2}{3k}\right\rfloor
= \left\lfloor \frac{m}{k} + \frac 2{3k}\right\rfloor
= \left\lfloor \frac {m}{k}\right\rfloor
$$
by being
$$
\frac{k-1}{k} + \frac {2}{3k} = \frac{3k-1}{3k} < 1;
$$
thus,
\begin{equation*}
\sum_{k=1}^{\lfloor\frac{n}3\rfloor} \left\lfloor \frac {n}{3k}\right\rfloor \mu(k)
= \sum_{k=1}^{m} \left\lfloor \frac {3m+2}{3k}\right\rfloor \mu(k)
= \sum_{k=1}^{m} \left\lfloor \frac {m}{k}\right\rfloor \mu(k) = 1.
\qedhere
\end{equation*}
\end{proof}

Then, we establish the following result, in which, as in Theorem~\ref{th:Mk},
we show an expansion of $M(n)$ as a sum of $\lfloor\frac{n}3\rfloor$ summands;
in every summand, only the M\"obius function and a coefficient
(related to the integer-part function) appear.
This will be a fruitful result, because, later in this paper, we will find
some alternative formulas and interesting properties for the coefficients.

\begin{theor}
\label{th:2M+3}
For $n \geq 3$, we have
\begin{equation}
\label{eq:2M+3}
2M(n)+3 = \sum_{k=1}^{\lfloor\frac{n}3\rfloor} g(n,k) \mu(k)
\end{equation}
with
\begin{equation}
\label{eq:g}
g(n,k) = 3 \left\lfloor \frac {n}{3k}\right\rfloor
- 2\left\lfloor \frac {n}{2k} - \frac 12 \right\rfloor.
\end{equation}
\end{theor}

\begin{proof}
Let us add $2$ times the expansion for $M(n)$ in Theorem~\ref{th:Mk}
plus $3$ times the expansion for $1$ in Proposition~\ref{th:1=3k}.
\end{proof}

The next result presents an alternative way for computing~$g(n,k)$:

\begin{propo}
\label{th:gnk}
For $k>0$ and $n\ge0$, let us take $n_0$ such that
$$
n\equiv n_0 \mod 6k,\quad 0\leq n_0<6k.
$$
Then
\begin{equation}
\label{eq:gnk}
g(n,k) = \left\{
\begin{array}{lll}
2 & \mbox{ if } & 0 \leq n_0 < k,\\
0 & \mbox{ if } & k \leq n_0 < 3k,\\
1 &  \mbox{ if } & 3k \leq n_0 < 5k,\\
-1 & \mbox{ if } & 5k \leq n_0 < 6k.
\end{array}
\right.
\end{equation}
\end{propo}

\begin{proof}
Le us decompose $n = n_0+6kn_1$, with  $0 \leq n_0 < 6k$. By~(\ref{eq:g}),
$$
g(n,k)
= 3 \left\lfloor \frac{n_0+6kn_1}{3k} \right\rfloor
- 2 \left\lfloor \frac{n_0+6kn_1-k}{2k} \right\rfloor.
$$
Then, it is clear that
\begin{align*}
\mbox{if\quad} & 0\leq n_0<k, & g(n,k)
&= 6n_1-2 \left\lfloor \tfrac{6k(n_1-1)}{2k}
+ \tfrac{5k+n_0}{2k} \right\rfloor = 2;\\
\mbox{if\quad} & k\leq n_0 < 3k, & g(n,k) &= 6n_1-6n_1 = 0;\\
\mbox{if\quad} & 3k \leq n_0 < 5k, & g(n,k) &= 3+6n_1-6n_1-2 = 1;\\
\mbox{if\quad} & 5k \leq n_0 < 6k, & g(n,k) &= 3+6n_1-6n_1-4 = -1.\qedhere
\end{align*}

\end{proof}

%%%%%%%%%%%%%%%%%%%%%%%%%%%%%%%%%%%%%%%%%%%%%%%
\section{Formulas in which both $M$ and $\mu$ appear}
\label{sec:Mmu}
%%%%%%%%%%%%%%%%%%%%%%%%%%%%%%%%%%%%%%%%%%%%%%%

Let us consider the function $g(n,k)$ for fixed $n$, i.e., as a function of~$k$.
In the following proposition, we show how $g(n,k)$ is constant when $k$ varies
a certain interval.

\begin{propo}
\label{th:gnkct}
Let $a$ and $n$ be positive integers, with $a<n$.
When $k$ varies in the interval
$$
\frac{n}{a+1} < k \leq \frac na
$$
the value of $g(n,k)$ remains constant. This value depends only
upon the remainder of $a$ modulus~$6$.
\end{propo}

\begin{proof}
Let us decompose $a = a_0+6a_1$ with $0 \leq a_0 < 6$.
If $\frac {n}{a+1}<k\leq \frac na$, then
$ka \leq n < k(a+1)$ and so
\begin{equation}
\label{eq:ineq1}
ka_0+6ka_1 \leq n < k(a_0+6a_1+1).
\end{equation}
Thus, $n = n_0+6ka_1$ for some $n_0$ verifying $0 \leq n_0 < 6k$.
By substituting this value of $n$ in~(\ref{eq:ineq1}), it becomes
\begin{equation}
\label{eq:ineq2}
ka_0 \leq n_0 < ka_0+k.
\end{equation}
By~(\ref{eq:gnk}), $g(n,k)$ takes the same value for all~$n_0$
that satisfies~(\ref{eq:ineq2}), and this value of $g(n,k)$ depends
only on~$a_0$.
\end{proof}

As a consequence of Proposition~\ref{th:gnkct}, we can define
the function
\begin{equation}
\label{eq:h=g}
h(a) = g(n,k) \quad \mbox{for}\quad \frac{n}{a+1} < k \leq \frac na.
\end{equation}
By using (\ref{eq:gnk}) (pay attention to $a_0$ in the proof
of Proposition~\ref{th:gnkct}), $h(a)$ takes these values:
\begin{equation}
\label{eq:h(a)}
h(a) = \left\{
\begin{array}{lll}
2, & \mbox{ if }& a\equiv 0 \mod 6, \\
0, & \mbox{ if }& a\equiv 1 \mod 6, \\
0, & \mbox{ if }& a\equiv 2 \mod 6, \\
1, & \mbox{ if }& a\equiv 3 \mod 6, \\
1, & \mbox{ if }& a\equiv 4 \mod 6, \\
-1, & \mbox{ if }& a\equiv 5 \mod 6.
\end{array}\right.
\end{equation}

Now, we will split the sum in (\ref{eq:2M+3}) in two parts,
introducing a parameter~$r$.
The first part will consist in the $\lfloor\frac{n}{r+1}\rfloor$
first terms of the sum in (\ref{eq:2M+3}). The second part will be,
of course, the summands that remain; they will be manipulated
in such way that we will get $r$ sumands in which only the functions
$h$ and $M$ appear. We can say that this is a \textit{mixed}
recursive formula for computing $M$: $M(n)$ is obtained from
$\mu(m)$ and $M(m)$ with $m<n$.

\begin{theor}
\label{th:2M+3gh}
Let $n$ and $r$ two integers satisfying $3\leq r\leq n-1$. Then
\begin{equation}
\label{eq:2M+3gh}
2M(n)+3 = \sum_{k=1}^{\lfloor\frac{n}{r+1}\rfloor} g(n,k)\mu(k)
+ \sum_{a=3}^r h(a)\left(M\left(\frac na\right)
- M\left(\frac n{a+1}\right)\right).
\end{equation}
\end{theor}

\begin{proof}
By~(\ref{eq:2M+3}) and~(\ref{eq:h=g}),
\begin{align*}
2M(n)+3 &= \sum_{k=1}^{\lfloor\frac{n}{3}\rfloor} g(n,k)\mu(k)
= \sum_{k=1}^{\lfloor\frac{n}{r+1}\rfloor} g(n,k)\mu(k)
+ \sum_{a=3}^r \sum_{k=\lfloor \frac {n}{a+1}\rfloor+1}^{\lfloor
\frac na\rfloor} g(n,k)\mu(k)
\\
&= \sum_{k=1}^{\lfloor\frac{n}{r+1}\rfloor} g(n,k)\mu(k)
+ \sum_{a=3}^r h(a)\sum_{k=\lfloor \frac{n}{a+1}
\rfloor+1}^{\lfloor \frac na\rfloor}\mu(k)
\\
&= \sum_{k=1}^{\lfloor\frac{n}{r+1}\rfloor} g(n,k)\mu(k)
+ \sum_{a=3}^r h(a)\left(M\left(\frac na\right)
- M\left(\frac n{a+1}\right)\right).
\qedhere
\end{align*}
\end{proof}

Prior to continue, let us note which would be the two limit
cases: of course, $r=2$ is Theorem~\ref{th:2M+3}; and,
when $r=n$, the sum indexed by $k$ in~(\ref{th:2M+3gh}) disappears.
Another particular case appears by taking $r = \lfloor\sqrt{n}\rfloor$;
thus~(\ref{eq:2M+3gh}) becomes
\begin{multline*}
M(n) = \frac12 \Biggl( -3 
+ \sum_{k=1}^{\lfloor\frac{n}{\lfloor\sqrt{n}\rfloor+1}\rfloor} g(n,k)\mu(k)
+ M\left(\frac n3\right) \\
+ \sum_{a=4}^{\lfloor\sqrt{n}\rfloor} (h(a)-h(a-1)) M\left(\frac na\right) 
- h(\lfloor\sqrt{n}\rfloor) M\left(\frac n{\lfloor\sqrt{n}\rfloor+1}\right) 
\Biggr), 
\end{multline*}
a formula that resembles~(\ref{eq:Stiel}) after isolating $M(n)$, 
but starting in $M(\frac n3)$ instead of~$M(\frac n2)$.

On the other hand, by splitting again the summand on the right
in~(\ref{eq:2M+3gh}), we can introduce a new parameter~$s$:

\begin{theor}
\label{th:2M+3sr}
Let $n$, $r$ and $s$ be three integers such that $s \geq 0$ and
$6s+9\leq r\leq n-1$.
Then
\begin{multline*}
%\label{eq:2M+3sr}
2M(n)+3 = \sum_{k=1}^{\lfloor\frac{n}{r+1}\rfloor} g(n,k)\mu(k) \\
+ \sum_{b=0}^s \left( M\left( \frac{n}{3+6b}\right)
- 2M\left( \frac{n}{5+6b}\right)
+ 3M\left( \frac{n}{6+6b}\right)
- 2M\left( \frac{n}{7+6b}\right)\right) \\
+ \sum_{a=6s+9}^r h(a)\left(M\left(\frac na\right)
- M\left(\frac n{a+1}\right)\right).
\end{multline*}
\end{theor}

\begin{proof}
Let us expand the summand $\sum_{a=3}^r$ in~(\ref{eq:2M+3gh}). In this way,
\begin{multline*}
2M(n)+3 = \sum_{k=1}^{\lfloor\frac{n}{r+1}\rfloor} g(n,k)\mu(k)
+ h(3)M\left(\frac n3\right) - h(3)M\left(\frac n{4}\right)
\\
\null+ h(4)M\left(\frac n4\right) - h(4)M\left(\frac n{5}\right)
+ h(5)M\left(\frac n5\right) - h(5)M\left(\frac n{6}\right)
\\
\null+ h(6)M\left(\frac n6\right) - h(6)M\left(\frac n{7}\right)
+ h(7)M\left(\frac n7\right) - h(7)M\left(\frac n{8}\right)
\\
\null+ h(8)M\left(\frac n8\right) - h(8)M\left(\frac n{9}\right)
+ \cdots
\\
\null+ h(6s+8)M\left(\frac n{6s+8}\right)
- h(6s+8)M\left(\frac n{6s+9}\right)
\\
\null+ \sum_{a=6s+9}^r h(a)\left(M\left(\frac na\right)
- M\left(\frac n{a+1}\right)\right).
\end{multline*}
By applying the values of $h$ according~(\ref{eq:h(a)}),
the result follows.
\end{proof}

Thus, in Theorem~\ref{th:2M+3sr} we have presented a two-parametric family of
recurrence relation for computing an isolated value of $M(n)$.
They provide mixed ways to calculate $M(n)$ using,
in part, previously computed (and stored) values of $M(m)$ for a
certain values of $m$, and another part that must be explicitly computed.
Eventually, a suitable election of parameters $r$ and $s$ (that may depend on~$n$)
will allow to get efficient methods of running this algorithm in a computer;
at this point, it is clear that a careful implementation must be performed,
taking into account the machine to be used.
%
The idea to use this expansion is as follows: The terms in the first sum
$\sum_{k=1}^{\lfloor\frac{n}{r+1}\rfloor}$ are directly evaluated.
The second sum $\sum_{b=0}^s$ is computed by using previously computed
and stored values of~$M$. And the third sum
$$
\sum_{a=6s+9}^r h(a)\left(M\left(\frac na\right)
- M\left(\frac n{a+1}\right)\right)
= \sum_{a=6s+9}^r h(a) \sum_{k=\lfloor
\frac n{a+1}\rfloor +1}^{\lfloor\frac na\rfloor}\mu(k)
$$
can be computed by using both methods, according the size of~$n$, $r$ and~$s$.

Finally, let us note which are the limit cases of the identity
established by Theorem~\ref{th:2M+3sr}: $s=-1$ (Theorem~\ref{th:2M+3gh});
$s=-1$ and $r=2$ (Theorem~\ref{th:2M+3}); $r=n$ ($\sum_b$~disappears);
and $r=6s+8$ ($\sum_k$~disappears).

%%%%%%%%%%%%%%%%%%%%%%%%%%%%%%%%%%%%%%%%%%%%%
\section{The function $H(n,m)$. Periodicity}
\label{sec:H}
%%%%%%%%%%%%%%%%%%%%%%%%%%%%%%%%%%%%%%%%%%%%%

In the previous sections (see Theorems~\ref{th:2M+3},
\ref{th:2M+3gh} and~\ref{th:2M+3sr}),
we often obtain expressions with the form
$\sum_{k=1}^m g(n,k)\mu(k)$.
Thus, in this section, we define a new function $H(n,m)$
(for non-negative integers $n$ and positive integers $m$) by taking
$$
H(n,m) := \sum_{k=1}^m g(n,k)\mu(k),
$$
and we are going to study some of its properties.
Also, we will use the following notation:
$$
C_m := 6 \cdot \lcm \{1,2,\ldots,m\}.
$$

First, let us see that, when we fix $m$ in the second variable of $H$,
the function is periodic with period~$C_m$.

\begin{propo}
\label{th:H}
For every non-negative integer $t$, we have
$$
H(n + t C_m,m) = H(n,m).
$$
\end{propo}

\begin{proof}
By being $g(n+{6k}t, k) = g(n,k)$ for $k=1,2,\ldots, m$, we have
$g(n+C_mt,k) = g(n,k)$. Thus, the result follows.
\end{proof}

The following result gives the value of $H(n,m)$ as a function of~$M(m)$.

\begin{propo}
%\label{th:HC}
For every non-negative integer $t$, we have
\begin{align*}
H(0 + t C_m,m) &= 2M(m),\\
H(1 + t C_m,m) &= 2M(m)-2,\\
H(2 + t C_m,m) &= 2M(m),\\
H(n + t C_m,m) &= 2M(m)+3, \quad \mbox { if }\ 2<n\leq m .
\end{align*}
\end{propo}

\begin{proof}
By Proposition~\ref{th:H}, without loss of generality, we can suposse $t=0$.
Then, it is enough for computing $H(n,m)$ for $0\leq n\leq m$.
First, let us analyze the cases $n=0,1,2,3$. By applying~(\ref{eq:gnk}), we have
\begin{align*}
H(0,m) &= \sum_{k=1}^m g(0,k)\mu(k) = 2\sum_{k=1}^m \mu(k) = 2M(m),
\\
H(1,m) &= \sum_{k=1}^m g(1,k)\mu(k) = 0\cdot \mu(1)+2\sum_{k=2}^m \mu(k) = 2M(m)-2,
\\
H(2,m) &= \sum_{k=1}^m g(2,k)\mu(k)
= 0\cdot \mu(1) + 0\cdot\mu(2) + 2\sum_{k=3}^m \mu(k) = 2M(m).
% \\ % Unnecessary
% H(3,m) &= \sum_{k=4}^m g(3,k)\mu(k)
% = 1\cdot \mu(1) + 0\cdot \mu(2) + 0\cdot\mu(3) + 2\sum_{k=4}^m \mu(k)
% = 2M(m)+3.
\end{align*}

For $n$ verifying $2 < n\leq m$, let us decompose
\begin{equation}
\label{eq:3sums}
H(n,m) = \sum_{k=1}^{\lfloor\frac n3\rfloor} g(n,k)\mu(k)
+ \sum_{k=\left\lfloor\frac n3\right\rfloor+1}^n g(n,k)\mu(k)
+ \sum_{k=n+1}^m g(n,k)\mu(k).
\end{equation}
Now, in the first sum, let us apply~(\ref{eq:2M+3});
in the second sum, let us use that, for
$\left\lfloor\frac n3\right\rfloor < k\leq n$ (i.e., $k\leq n<3k$)
we have $g(n,k)=0$ (see~(\ref{eq:gnk})); and, finally,
for the third sum, let us note that $g(n,k)=2$ for $0\leq n<n+1\leq k$.
In this way,~(\ref{eq:3sums}) becomes
\begin{equation*}
H(n,m) = (2M(n)+3) + 0 + 2\left(M(m)-M(n)\right) = 2M(m)+3.
\qedhere
\end{equation*}
\end{proof}

%%%%%%%%%%%%%%%%%%%%%%%%%%%%%%%%%%%%%%%%%%%%%%%%%%
\begin{thebibliography}{15}
%%%%%%%%%%%%%%%%%%%%%%%%%%%%%%%%%%%%%%%%%%%%%%%%%%

% References with style "Vancouver", as requested by The Open Mathematics Journal.
% The references are ordered according the order of its citation in the paper.

\bibitem{Ap}
{Apostol TM}.
{Introduction to Analytic Number Theory}.
New York: Springer-Verlag; 1976.

\bibitem{Mer}
{Mertens F}.
{\"Uber eine zahlentheoretische Funktion}.
{Sitzungsberichte Akad. Wiss. Wien IIa}. 
1897; {106}: 761--830.

\bibitem{Her}
{Hermite C}. 
{Correspondence d'Hermite et de Stieltjes: publi\'ee par le soins de B.~Baillaud et H.~Bourget}. Vol.~1. 
Paris: Gauthier-Villars; 1905. 
%% Also available in:
% The Cornell University Library Digital Collections, 2005,
% http://name.umdl.umich.edu/AAN9223.0001.001

\bibitem{Lit}
{Littlewood JE},
{Quelques cons\'equences de l'hypoth\`ese que la fonction $\zeta(s)$
n'a pas de z\'eros dans le demi-plan $\operatorname{Re}(s)>\frac 12$}.
{C. R. Acad. Sci. Paris}. 
1912; {154}: 263--266.

\bibitem{Tit}
{Titchmarsh EC}.
{The Theory of the Riemann Zeta-Function}.
Oxford: At the Clarendon Press; 1967.
%% Existe una segunda edici\'on (con prefacio de D. R. Heath-Brown),
%% pero es de 1986.

\bibitem{OdRie}
{Odlyzko AM, te Riele H}.
{Disproof of the Mertens conjecture}.
{J. Reine Angew. Math.} 
1985; {357}: 138--160.

\bibitem{Pi}
{Pintz J}.
{An effective disproof of the Mertens conjecture}.
In: {Journ\'ees Arithm\'etiques} (Besan\c{c}on, France, 1985).
{Ast\'erisque}. 
1987; {147-148}: 325--333, 346.

\bibitem{KoRie}
{Kotnik T, te~Riele H}. 
{The Mertens conjecture revisited}. 
In: {Algorithmic number theory}.
{Lecture Notes in Comput. Sci.} {4076}. 
Berlin: Springer; 2006. p.~156--167.
% DOI 10.1007/11792086_12. 
% http://dx.doi.org/10.1007/11792086_12

\bibitem{Ma}
{el Marraki M}.  
{Fonction sommatoire de la fonction de M\"obius, 3. Majorations asymptotiques effectives fortes}.
{J. Th\'eor. Nombres Bordeaux}. % {Journal de Th\'eorie des Nombres de Bordeaux} 
1995; {7}: 407--433.
% http://archive.numdam.org/article/JTNB_1995__7_2_407_0.pdf

\bibitem{KoLu}
{Kotnik T, van de Lune J}.
{On the order of the Mertens function}. 
{Experiment. Math.}
2004; {13}: 473--481. 
% http://www.expmath.org/expmath/volumes/13/13.4/Kotnik.pdf

\bibitem{Lan}
{Landau E}. 
{Handbuch der Lehre von der Verteilung der Primzahlen}. 
Leipzig \& Berlin: Teubner; 1909.

\bibitem{DeRi}
{Del\'eglise M, Rivat J}.
{Computing the summation of the M\"obius function}.
{Experiment. Math.} % Experimental Mathematics 
1996; {5}: 291--295.
% http://www.expmath.org/restricted/5/5.4/rivat.ps

\bibitem{Neu}
{Neubauer G}. 
{Eine empirische Untersuchung zur Mertensschen Funktion}. 
{Numer. Math.} % Numerische Mathematik 
1963; {5}: 1--13. 
% DOI 10.1007/BF01385874. 
% http://dx.doi.org/10.1007/BF01385874

\bibitem{Dr}
{Dress F}. 
{Fonction sommatiore de la fonction de M\"obius. I. Majorations exp\'erimentales}. 
{Experiment. Math.} 
1993; {2}: 89--98. 
% http://www.expmath.org/restricted/2/2.2/dress.ps

\bibitem{Rie}
{te~Riele H}. 
{On the history of the function $M(x)/\sqrt{x}$ since Stieltjes}. 
In: Thomas Jan Stieltjes: {\OE}uvres Compl\`etes/Collected Papers. Vol.~1. 
Berlin: Springer-Verlag; 1993. p. 69--79.

\end{thebibliography}

\thispagestyle{final}

%%%%%%%%%%%%%%%%%%%%%%%%%%%%%%%%%%%%%%%%%%%%%%%%%%
\end{document}
%%%%%%%%%%%%%%%%%%%%%%%%%%%%%%%%%%%%%%%%%%%%%%%%%%


